
\documentclass[11pt]{article}
\usepackage[T1]{fontenc}
\usepackage[utf8]{inputenc}
\usepackage{lmodern}
\usepackage{geometry}
\usepackage{microtype}
\usepackage{amsmath,amsthm,amssymb,amsfonts,mathtools}
\usepackage{enumitem}
\usepackage[colorlinks=true,linkcolor=blue,citecolor=teal,urlcolor=blue]{hyperref}
\geometry{margin=1in}
\title{\textbf{Arithmetic of Platonic Surfaces}\\[2pt]\large Algorithms, Fields, and a Planar-Points Pipeline}
\author{}
\date{}
\newtheorem{theorem}{Theorem}[section]
\newtheorem{proposition}[theorem]{Proposition}
\newtheorem{lemma}[theorem]{Lemma}
\theoremstyle{definition}
\newtheorem{definition}[theorem]{Definition}
\theoremstyle{remark}
\newtheorem{remark}[theorem]{Remark}
\newcommand{\Hh}{\mathbb{H}}
\newcommand{\Z}{\mathbb{Z}}
\newcommand{\Q}{\mathbb{Q}}
\newcommand{\PSL}{\mathrm{PSL}}
\newcommand{\tr}{\mathrm{tr}}
\begin{document}
\maketitle
\begin{abstract}
We revisit arithmeticity for surface groups arising from ideal polyhedra, record a computation-friendly Penner--Rivin algorithm, explain how to compute the invariant trace field from shears, and add a practical pipeline that starts from a planar point set and ends with a uniformizing Fuchsian group. A worked ``off-center hexagon'' example closes the loop (shears, field, volume).
\end{abstract}
\section{Trace fields and arithmeticity}
For a finitely generated Fuchsian group $\Gamma<\PSL_2(\mathbb{R})$ and a lift $\widetilde\Gamma<\mathrm{SL}_2(\mathbb{R})$, the \emph{invariant trace field} is
\begin{equation*}
k(\Gamma)=\Q\big(\{\tr(\tilde g^2):\tilde g\in\widetilde\Gamma\}\big).
\end{equation*}
It is independent of the lift and stable under commensurability. In the \emph{cusped} case, arithmetic $\iff k(\Gamma)=\Q$ (so $\Gamma$ is commensurable with $\PSL_2(\Z)$). From an ideal triangulation with shear coordinates $Z_e$, the Penner--Rivin product expresses $\tr(g)$ as a Laurent polynomial in $e^{Z_e/2}$, so $k(\Gamma)\subseteq \Q(e^{Z_e/2})$.
\section{Planar points $\to$ convex hull $\to$ group}
Given $P=\{z_i\}\subset\C$:
\begin{enumerate}[leftmargin=2em]
\item \textbf{Back-project} to $S^2$ by inverse stereographic projection; take the Euclidean \textbf{convex hull}.
\item \textbf{Triangulate faces} by a fan; each oriented edge $x\to y$ with adjacent vertices $a$ (left) and $b$ (right) carries the \emph{complex shear--bend}
\begin{equation*}
\tau_{xy}=\log\!\frac{(a-x)(b-y)}{(a-y)(b-x)}.
\end{equation*}
Its real part is the Penner shear; the imaginary part is the exterior dihedral (zero for face diagonals).
\item Feeding the shears $Z_e=\Re\tau_e$ into the Penner--Rivin algorithm gives a generating set of matrices for the holonomy $\Gamma<\PSL_2(\mathbb{R})$.
\end{enumerate}
\section{Worked example: the off-center hexagon}
Let $\omega=e^{\pi i/3}$ and $z_i=1+\omega^i$ ($i=0,\dots,5$). Add the north pole to obtain a hexagonal pyramid. With a base fan triangulation, lateral edge shears vanish while equatorial shears take two nonzero values (up to sign), yielding a nonarithmetic group with
\begin{equation*}
k(\Gamma)=\Q(\sqrt{3}).
\end{equation*}
The volume is the sum of four Bloch--Wigner values for the shapes $(z_0-z_{k+1})/(z_k-z_{k+1})$, $k=1,2,3,4$.
\paragraph{Code.} The repository provides a python module that implements this pipeline end-to-end (convex hull, shears, holonomy, field signature, volume, mesh, and a PNG renderer with lights and soap-film iridescence).
\end{document}
